\documentclass[a4paper, 12pt, french]{article}

% ========================
% IMPORTATION DES PACKAGES
% ========================

\usepackage[T1]{fontenc}
\usepackage[utf8]{inputenc}
\usepackage[margin=2cm]{geometry}
\usepackage{lmodern}
\usepackage{theorem}
\usepackage{amsfonts, amsmath, amssymb}
\usepackage{babel}

\pagestyle{headings}

\renewcommand{\familydefault}{\sfdefault}

\theoremstyle{break}
\theoremheaderfont{\scshape}
\theorembodyfont{\upshape}

\newtheorem{enon}{Énoncé}[section]
\newtheorem{preuve}{Preuve}[section]

\title{MÉSURE ET INTÉGRATION}
\author{}
\date{\today}

\begin{document}
\maketitle

\begin{center}\section{Théorie des ensembles}\end{center}
\begin{enon}
       Soit $(A_{n})_{n \in \mathbb{N}} \subseteq P(X)$.
       \begin{enumerate}
               \item Montrer que  $\displaystyle\bigcup_{n \in \mathbb{N}} A_{n} =  \bigsqcup_{n \in \mathbb{N}} B_{n}$.
               \item Est ce qu'on peut écrire  $\displaystyle\bigcup_{n \in \mathbb{N} } A_{n}$ comme une suite croissante ? respectivement comme une suite décroissante ?
       \end{enumerate}
\end{enon}
\begin{preuve}
      \begin{enumerate}
              \item Soit $(A_{n})_{n \in \mathbb{N}} \subseteq P(X)$.\\ Posons $ B_{0}= A_{0}, \; B_{n} = A_{n} \backslash \bigcup_{n=1}^{n-1} A_{n}$ pour $n \geq 1$. On peut voir facilement que les $B_{n}$ sont deux à deux disjoints.\\
              D'abord montrons que $\displaystyle\bigcup_{n \in \mathbb{N}} A_{n} \subseteq \bigsqcup_{n \in \mathbb{N}} B_{n}$.\\
              Soit $ x \in \displaystyle\bigcup_{n \in \mathbb{N}} A_{n}$, alors il existe $n_{0}$ tel que $x \in A_{n_{0}}$. Alors $x \not\in A_{k<n_{0}}$, par suite $x \not\in \displaystyle\bigcup_{k<n_{0}} A_{k}$. Donc $x \in A_{n_{0}} \backslash\displaystyle\bigcup_{n \in \mathbb{N}} A_{n} = B_{n_{0}}$.
              Ainsi $x \in \displaystyle\bigsqcup_{n \in \mathbb{N}} B_{n}$.
      \end{enumerate}
\end{preuve}
\end{document}